\section{Discussion}\label{sec:disc}

\subsection{What the results permit}\label{sec:disc:uses}

The theorems concern structured nonconvex instances under the specific finite
laws constructed before sampling. Independence is in the coordinates of the
perturbation model: original coefficients for ambient noise, core
coefficients for core-only noise, and factor coefficients for aligned noise.
Within these contracts they give three kinds of guarantees.

First, they return exact global solutions of sampled instances, with expected
bit work bounded by the structural parameters and displayed numerical
ratios. If uncertain or estimated costs follow the actual law and hypotheses
of a theorem, the sampled instance can itself be the problem of interest.
An exogenous cost distribution is covered only when it matches that
contract; arbitrary uncertainty distributions do not inherit these work
bounds.

Second, for routes valid at every positive noise scale,
\cref{prop:model:regret,sec:model:original} give an exactly feasible point and
an original-objective certificate of any prescribed accuracy $\varepsilon$.
The recipe uses widths in the perturbation coordinates and the exact
Gaussian-like support factor $b+20$. The expected work is polynomial in
$1/\varepsilon$ for fixed structure and fixed other numerical parameters;
it is polynomial in $I$ and $1/\varepsilon$ only when those numerical
parameters are polynomially bounded. Implicit graphs use exact algebraic
lifts of rational free points. The strong-field theorem retains the regret
bound, but its lower noise-scale condition prevents this calibration in
general.

Third, the proofs give certificates checked exactly and sound on every draw:
corrected-corner bounds, critical regions of parametric convex programs,
strongly convex patches after sound coordinate fixing, excluded-region and
competing-label gaps, and optimal-face certificates for integer flow and TU
recourse. Probability bounds control when these certificates pass. They
could serve as termination tests inside other global solvers after branching
has localized the optimum; their practical effectiveness is unestablished.

The structures have natural application interpretations. A small core with
convex, network-flow, or TU recourse describes two-stage problems in which a
few nonconvex first-stage decisions enter residual costs while residual
feasibility stays fixed. General coupling to second-stage constraints is
excluded. Separable convex costs plus a low-rank concave term can describe
economies of scale in aggregate quantities with many integer variables.
Orders with binary variables describe isotonic fitting with switching;
simplex products describe allocation; graph equalities describe actuator
models. Each application must meet the actual law, curvature, domain, and
oracle hypotheses. These examples provide no performance evidence.

Core-only noise locates the core, while a residual modulus or an optimal-face
certificate supplies point output in our theorems. These are sufficient
structures, rather than necessary conditions for every problem. The
exact-arithmetic companion also gives full selected-point oracles for
residual-convex cubics without a modulus, and supplied-convexifier point
oracles under its stated degree and global-convexity
hypotheses~\citep{companion-exact-arithmetic}. The general arithmetic
obstruction of \cref{thm:lim:posslp} starts at degree four.

\subsection{Limitations}\label{sec:disc:limits}

Numerical ratios such as $L/\sigma$ and $\nu\diam(X)/\sigma$ enter through
their magnitudes. The FPT claims parameterize them jointly with structure;
for a fixed bounded ratio they are FPT in structure. A polynomially growing
ratio in $(1+R)^{O(k)}\poly(I)$ gives $I^{O(k)}$, rather than FPT in $k$
alone. The sparse bounds are polynomial for each fixed width, with a
dimension factor necessary for the specified retention rule
(\cref{thm:lim:width}).

Finite laws need enough resolution for the base-only format, stopping-depth,
and fallback budgets. Polynomial-precision routes admit one resolution per
input length, scaled by the supplied noise scale (\cref{prop:model:uniform}).
General nonlinear core-only flow and TU recourse retain a budget depending
on a supplied core-size bound. Arbitrary coarse laws are not covered. Most
closure routes handle ties with a rare same-draw fallback; lattice and
strong-field methods finish directly.

Several premises are supplied: sharper curvature bounds, cores, tree
decompositions, convexifiers, residual convexity or moduli, and recourse
oracles valid under every required restriction. Certificate encoding length
belongs in $I$ and verification work in the runtime. A general polynomial
inequality test is not hidden in the input model.

Exact output also has a representation contract. Ordinary implicit or
charted descriptors are compact and support evaluation to any precision;
they do not provide exact threshold comparisons or active-set decisions.
Fallback outputs have base-exponential bounds on every draw, and global
proof records can also be large; their expected lengths enter the work
bounds. Small-core algebraic outputs retain their per-draw parameter bounds.
Component outputs give common-root data for each component and a symbolic
sum of values; exact comparisons of unrelated algebraic sums are not
claimed. The companion instead gives Cauchy names of a canonical selected
point with one per-draw work factor controlling all precision queries under
one base-selected law. Our finite descriptor and global proof record are
an output distinction, rather than by themselves an improvement in
optimization time.

The constants are large. In particular, the nonquadratic strong-field
sufficient condition contains the conservative component-solver base
$c_d=2^{10000D}$, with $D$ the least even integer greater than $d$, in both
the derivative-variation and grid-size requirements. The quadratic case
uses $a_i=3$ for continuous coordinates and has a smaller sufficient scale.
These conditions define proved regimes, but their conservative size and the
absence of computational evidence leave practical performance unsettled.

\subsection{Open questions}\label{sec:disc:open}

\emph{Width.} Can exact sparse polynomial optimization under independent
linear noise have expected work FPT jointly in width and numerical ratios?
\Cref{thm:lim:width,prop:lim:local} rule out the analyzed global-error
retention rule and its naive bag-local replacement. Certified conditional
values are one sufficient way to remove the ambient dimension from this
count (\cref{prop:sp:conditional}); the current grid min-marginals do not
supply them. They are not a proved necessary ingredient of every FPT
algorithm. The deterministic growth-conditioned graded-grid companion gives
another relevant starting point~\citep{companion-decomposition-aware}.
A linear small-growth tail alone does not integrate a generic
inverse-growth work factor $\kappa^r$, $r>1$, to an FPT expectation without
a cutoff-dependent factor: its truncated-moment estimate retains
$K^{1-1/r}$. This limits that integration argument, rather than other
algorithms (\cref{sec:sp:width,rem:qp:moments}).

\emph{Uniform ambient noise.} Can the low-rank route be FPT jointly in
negative inertia and numerical ratios under independent uniform ambient
noise? \Cref{thm:lim:ambient} excludes a dimension-free bound for the
specified auxiliary counts; it does not exclude another algorithm.

\emph{Constraints.} Which sparse linear constraints beyond the proved graph,
simplex, and order domains admit exact smoothed algorithms? The affine
hardness construction needs restrictions beyond sparsity
(\cref{thm:lim:constraints}), and the TU diagonal example calls for counts
adapted to feasible faces (\cref{ex:lim:coupled}).

\emph{Residual points.} Which structures beyond the positive cubic,
convexifier, modulus, and integer-face results permit efficient residual
point evaluation under core-only noise? Convex quartics already carry the
Square Root Sum and $\PosSLP$ implications of \cref{thm:lim:posslp}.
A continuous optimal-face certificate is one possible direction.

\emph{Formats and precision.} Which further classes admit short common-root
algebraic outputs on every draw? Can general nonlinear core-only flow and TU
recourse use polynomial sampling precision while retaining the current
parameter dependence, as bilinear coupling does?

\emph{Computation.} Implementations with certified convex solvers, compact
proof records, and smaller constants are needed to assess the certificates
as practical termination tests.
